%% kapa: exam=1 round=29 year=2018 period=2 form=A subject=appraisal_law subject_ko=감정평가관계법규 seq=12 number=12
%% source: Q-Net 감정평가사 기출문제 (cst00309) artlSeq=5202808 file=extracted/1008/5202808/2018년 제29회 감정평가사 1차시험 기출문제/2018년 제29회 감정평가사 1차시험 기출문제/2교시 A형.pdf
%% conversion: kordoc 4.24.51 -> tools/gen1.py (auto); review=proofread-sonnet; flags=none
\begin{question}{12}
국토의 계획 및 이용에 관한 법령상 시범도시에 관한 설명으로 옳은 것은?
\begin{choices}
  \item 국토교통부장관과 시ㆍ도지사는 시장ㆍ군수ㆍ구청장의 신청을 받아 시범도시를 지정할 수 있다.
  \item 시범도시사업의 시행을 위하여 필요한 경우에는 시범도시사업의 예산집행에 관한 사항을 도시ㆍ군계획조례로 정할 수 있다.
  \item 시범도시를 공모할 경우 이에 응모할 수 있는 자는 특별시장ㆍ광역시장ㆍ특별자치시장ㆍ특별자치도지사ㆍ시장ㆍ군수ㆍ구청장 또는 주민자치회이다.
  \item 국토교통부장관은 시범도시를 지정하려면 설문조사ㆍ열람 등을 통하여 주민의 의견을 들은 후 관계 지방자치단체장의 의견을 들어야 한다.
  \item 국토교통부장관은 시범도시사업계획의 수립에 소요되는 비용의 전부에 대하여 보조 또는 융자할 수 있다.
\end{choices}
\end{question}
\answer{2}
